\hmsubsection{Risk Management}{FS}{}\label{sec:risk-management}
This section presents the risks identified at the start of the project
and the planned mitigation measures. A 3 by 3 matrix is used instead of
a 5 by 5 matrix. A smaller matrix is easier to apply consistently across
the team, and the supervisor recommended it for projects of this size.
With three levels each for likelihood and consequence, the risk index is
between 1 and 9, which maps to three action classes: low (1--2), medium
(3--4), and high (6--9).
\hmsubsubsection{How we rate the risks}{FS}{}
The risk assessment follows the HAZID approach from IEC 61882
\cite{hazid}. Each risk is assigned two ratings on a 1--3 scale.
Likelihood (L) represents the estimated probability that the risk will
occur during the project. Consequence (C) represents the severity of the
impact if it occurs. The risk index is the product of the two:
\[
\text{RI} = L \times C, \qquad L, C \in \{1, 2, 3\}, \qquad \text{RI} \in [1, 9]
\]
The three levels mean the same thing on both axes:
\begin{itemize}
\item \textbf{1 -- low.} Unlikely (L) or minor (C). For likelihood,
the event is not expected to occur during the semester. For consequence,
the issue can be resolved in one or two days without loss of sprint
progress.
\item \textbf{2 -- medium.} Possible (L) or noticeable (C). For
likelihood, the event has been observed on similar projects, including
Sortify \cite{sortify}. For consequence, recovery requires between a few
days and one sprint.
\item \textbf{3 -- high.} Likely (L) or serious (C). For likelihood,
the event is expected to occur unless it is actively prevented. For
consequence, the event threatens the delivery deadline or makes the
system unusable.
\end{itemize}
The 3 by 3 grid was selected instead of the 5 by 5 grid for two reasons.
First, the difference between adjacent levels is large enough to support
consistent team ratings. With a 5 by 5 grid, the additional granularity
led to discussion about adjacent numerical ratings rather than
mitigation actions. Second, the 3 by 3 grid produces three action
classes, which map clearly to sprint planning. A risk is either monitored
for now (low), mitigated in the next sprint (medium), or addressed in
the current week (high). The matrix is shown in Table
\ref{tab:risk-matrix}.
\begin{table}[h]
\centering
\small
\begin{tabular}{c|ccc}
\toprule
& \textbf{C = 1} & \textbf{C = 2} & \textbf{C = 3} \\
& (minor) & (noticeable) & (serious) \\
\midrule
\textbf{L = 3 (likely)} & 3 (medium) & 6 (high) & 9 (high) \\
\textbf{L = 2 (possible)} & 2 (low) & 4 (medium) & 6 (high) \\
\textbf{L = 1 (unlikely)} & 1 (low) & 2 (low) & 3 (medium) \\
\bottomrule
\end{tabular}
\caption{The 3 by 3 risk matrix used in this project. Rows are
likelihood, columns are consequence, cells show the risk index and its
action class. Low means that the risk is monitored, medium means that it
is mitigated within the next sprint, and high means that it is addressed
in the current week.}
\label{tab:risk-matrix}
\end{table}
\subsubsection*{When we review the register}
The risk register is reviewed at the end of every sprint during the
retrospective. New risks are added at the bottom of the table. Risks
that are no longer relevant are retained in the table with a note, so
the decision history remains traceable. This is important for the report
because it allows design decisions to be linked to the sprint in which a
risk was rated and the sprint in which mitigation was performed.
Three risks changed rating during the project. R-02 (inverse kinematics
accuracy) started at L=3, C=3, RI=9 because Sortify reported this as its
main limitation. After the closed-form solver and sag compensation were
implemented, the likelihood was reduced to 2 and the index to 6. R-04
(elbow motor overheats at the scan pose) was not included in the
original register. It was added after the motor shut down during a long
demonstration. R-10 (MobileNetV2 inference time) was retired in week 6
after measurements of 200--500 ms per frame on the Raspberry Pi 5 led to
the selection of the classical pipeline.

\subsubsection*{What each risk means in practice}
R-01 (schedule slip) is a general project-management risk. It was rated
3 by 3 because the team consists of five members across two disciplines,
the assignment is open-ended, and the semester is short. The weekly
cadence of client and supervisor meetings is the primary mitigation.
Sprint retrospectives provide an additional opportunity to identify
schedule drift early.
R-02 (inverse kinematics accuracy) was the highest-rated technical
risk at the start of the project. Sortify \cite{sortify} reported their
arm hitting target points with visible deviation, and they could not
complete fully autonomous pick-and-place because of it. A closed-form
geometric solver was selected instead of an iterative one to reduce the
likelihood of the same issue. After the solver and sag model were in
place, the gripper reached commanded points within roughly 1 cm on the
workspace, which satisfies the requirement for a 50 mm ball. The risk is
still retained in the table because calibration can drift if the arm is
bumped or remounted.
R-03 (lighting) is the most common reason colour-based detection
fails. Standard HSV thresholding only works inside a narrow range of
lighting, and shadows from the arm itself can cause detection failure.
CLAHE on every frame mitigates most of this effect. Multiple HSV ranges
for the shadow case address the remaining variation. This was tested
between 300 and 700 lux on the bench and remained functional.
R-04 (elbow overheating) emerged during integration. The arm spends most
of its idle time at the scan pose. In the initial pose, the elbow motor
had to hold the forearm and the camera against gravity, and long
stationary demonstrations could cause the motor to latch an overheating
error and stop. A 300~mA current limit was tested at the scan pose, but
it did not reduce the thermal load sufficiently and was therefore not
retained in the final system.
Reducing the mass of the outer arm components improved the situation
somewhat, but the main mitigation for idle scan-pose heating was to
revise the scan pose so that the L\textsubscript{2} structure rests on
the L\textsubscript{1} piece. This mechanically supported posture reduces
the holding torque required from M3 while preserving the camera's view of
the balls. The system was able to sort 10 balls consecutively, but the
motors still became somewhat hot during movement. Operational monitoring
therefore remains relevant because the scan-pose change mitigates
stationary heating rather than eliminating movement-related thermal risk,
and no instrumented long-duration thermal characterisation was performed.
R-05 (serial desync) is mitigated primarily by the protocol design.
Newline-terminated lines are used, so if a byte is lost, the next
newline resynchronises the parser. The host adds a timeout-and-retry
mechanism. The available lab observations did not indicate
desynchronisation; the risk is retained mainly as a guard against partial
USB cable disconnection.
R-06 (motor failure) and R-07 (power supply) are hardware risks rated
low because spare parts and a regulated bench supply are available. They
remain in the table because losing a motor mid-sprint would delay the
project by several days. The factory-reset procedure in
Section \ref{sec:motor-config} also functions as a mitigation: two
motors arrived with latched overload errors and both were recovered.
R-08 (data loss) is a low risk because of Git. It is included because
calibration data is as important as the source code, and at the start of
the project the calibration was stored locally on one laptop. Once it
was moved into JSON files in the repository, this risk dropped to L=1.
R-09 (calibration drift) was the most common cause of a missed pick
in the first integration tests. The touch calibration procedure can
be re-run in 5--10 minutes, so drift no longer ends a test session. The
system is recalibrated at the start of each test day as a precaution.
R-10 (MobileNetV2) is the only retired risk. It is retained in the table
because the decision to remove MobileNetV2 was one of the more
significant early design choices and should remain traceable. The
section on similar systems (Section \ref{sec:related-systems}) and the
related-work section (Section \ref{sec:related-detection}) both refer
back to this entry.
\subsubsection*{Summary count}
Across the ten risks in the register, four are high (RI = 6 or above),
five are medium (RI = 3 or 4), and one is low (RI = 2). None of the
risks are at the top of the matrix (RI = 9) at the current state of
the project. R-01 (schedule slip) is the closest, and the remaining
high-rated risks already have mitigations in place.
